% GENERATED FILE. Edit canonical lessons, then run npm run generate:books.
% canonical-source-hash: fnv1a64:68f155b98f932e4a
% canonical-lessons: MR-C310-vrttapatra, MR-C310-dainik, MR-C310-saptahik, MR-C310-mathla, MR-C310-chapne

\chapter{The Morning Paper}
\label{ch:mr-the-morning-paper}
\hlchaptermodality{310}

\begin{chapteropening}
\textbf{By the end of this chapter:} \emph{I can talk about newspapers, headlines and printing.}

Complete the last lesson of chapter 310: I can talk about newspapers, headlines and printing.
\end{chapteropening}

\section[\texorpdfstring{\mr{वृत्तपत्र} (vṛttapatra)}{वृत्तपत्र (vṛttapatra)}]{\mr{वृत्तपत्र} (vṛttapatra) — a newspaper}

\label{lesson:MR-C310-vrttapatra}



\begin{quote}
Before the new one: say the Marathi for to pass away, then the Marathi for to arise.
\end{quote}

\subsection*{You'll want to know: \mr{वृत्तपत्र}}
\textbf{\mr{वृत्तपत्र}} — \emph{vṛttapatra} — \textquotedblleft{}a newspaper\textquotedblright{}.

\begin{quote}
\mr{आजोबा} \mr{रोज} \mr{सकाळी} \mr{वृत्तपत्र} \mr{वाचतात}.
\end{quote}

\begin{quote}
\emph{ājobā roj sakāḷī vṛttapatra vāctāt.}
\end{quote}

\begin{quote}
Grandpa reads the newspaper every morning.
\end{quote}

Neuter.

Also \mr{वर्तमानपत्र}, and in speech \mr{पेपर}.

\begin{cousinweb}
\textbf{\mr{वृत्तपत्र}} (\emph{vṛttapatra}) is a modern compound built from Sanskrit \emph{vṛtta}, 'what has happened, an event', and \emph{patra}, 'a leaf, a letter' (Monier-Williams).
\end{cousinweb}

\begin{grammarlens}[title={What you've built}]
One more word for telling what happened.
\end{grammarlens}

\subsection*{Your turn}
\begin{itemize}
\raggedright
  \item \emph{Say it:} \emph{vṛttapatra}
  \item \emph{Say it:} \emph{vṛttapatra}, once more
  \item \emph{Say it:} \emph{ājobā roj sakāḷī vṛttapatra vāctāt.} — Grandpa reads the newspaper every morning.
  \item \emph{Recall:} say the Marathi for to pass away, then the Marathi for to arise, then say \emph{vṛttapatra} again
\end{itemize}

\begin{tcolorbox}[breakable,colback=teal!4,colframe=teal!35!black,title={Before you move on}]
What does \mr{वृत्तपत्र} mean? (\textquotedblleft{}A newspaper\textquotedblright{}.) Say it once more.
\end{tcolorbox}
\section[\texorpdfstring{\mr{दैनिक} (dainik)}{दैनिक (dainik)}]{\mr{दैनिक} (dainik) — a daily paper; daily}

\label{lesson:MR-C310-dainik}



\begin{quote}
Before the new one: say the Marathi for to arise, then the Marathi for a newspaper.
\end{quote}

\subsection*{You'll want to know: \mr{दैनिक}}
\textbf{\mr{दैनिक}} — \emph{dainik} — \textquotedblleft{}a daily paper; daily\textquotedblright{}.

\begin{quote}
\mr{माझे} \mr{बाबा} \mr{रोज} \mr{दोन} \mr{दैनिकं} \mr{वाचतात}.
\end{quote}

\begin{quote}
\emph{mājhe bābā roj don dainikaṁ vāctāt.}
\end{quote}

\begin{quote}
My father reads two daily papers every day.
\end{quote}

Neuter.

Neuter as a noun, plural \mr{दैनिकं}; as an adjective: \mr{दैनिक} \mr{काम}, daily work.

\begin{cousinweb}
\textbf{\mr{दैनिक}} (\emph{dainik}) is a tatsama, taken from Sanskrit \emph{dainika}, 'daily' (Monier-Williams), built on \emph{dina}, 'a day'.
\end{cousinweb}

\begin{grammarlens}[title={What you've built}]
One more word for telling what happened.
\end{grammarlens}

\subsection*{Your turn}
\begin{itemize}
\raggedright
  \item \emph{Say it:} \emph{dainik}
  \item \emph{Say it:} \emph{dainik}, once more
  \item \emph{Say it:} \emph{mājhe bābā roj don dainikaṁ vāctāt.} — My father reads two daily papers every day.
  \item \emph{Recall:} say the Marathi for to arise, then the Marathi for a newspaper, then say \emph{dainik} again
\end{itemize}

\begin{tcolorbox}[breakable,colback=teal!4,colframe=teal!35!black,title={Before you move on}]
What does \mr{दैनिक} mean? (\textquotedblleft{}A daily paper; daily\textquotedblright{}.) Say it once more.
\end{tcolorbox}
\section[\texorpdfstring{\mr{साप्ताहिक} (sāptāhik)}{साप्ताहिक (sāptāhik)}]{\mr{साप्ताहिक} (sāptāhik) — a weekly}

\label{lesson:MR-C310-saptahik}



\begin{quote}
Before the new one: say the Marathi for a newspaper, then the Marathi for a daily paper.
\end{quote}

\subsection*{You'll want to know: \mr{साप्ताहिक}}
\textbf{\mr{साप्ताहिक}} — \emph{sāptāhik} — \textquotedblleft{}a weekly\textquotedblright{}.

\begin{quote}
\mr{हे} \mr{साप्ताहिक} \mr{दर} \mr{गुरुवारी} \mr{येतं}.
\end{quote}

\begin{quote}
\emph{he sāptāhik dar guruvārī yetaṁ.}
\end{quote}

\begin{quote}
This weekly comes out every Thursday.
\end{quote}

Neuter.

Built on \mr{सप्ताह}, a week; compare \mr{मासिक}, a monthly.

\begin{cousinweb}
\textbf{\mr{साप्ताहिक}} (\emph{sāptāhik}) is built on Sanskrit \emph{saptāha}, 'a week' (Monier-Williams), with the adjective ending \emph{-ika}.
\end{cousinweb}

\begin{grammarlens}[title={What you've built}]
One more word for telling what happened.
\end{grammarlens}

\subsection*{Your turn}
\begin{itemize}
\raggedright
  \item \emph{Say it:} \emph{sāptāhik}
  \item \emph{Say it:} \emph{sāptāhik}, once more
  \item \emph{Say it:} \emph{he sāptāhik dar guruvārī yetaṁ.} — This weekly comes out every Thursday.
  \item \emph{Recall:} say the Marathi for a newspaper, then the Marathi for a daily paper, then say \emph{sāptāhik} again
\end{itemize}

\begin{tcolorbox}[breakable,colback=teal!4,colframe=teal!35!black,title={Before you move on}]
What does \mr{साप्ताहिक} mean? (\textquotedblleft{}A weekly\textquotedblright{}.) Say it once more.
\end{tcolorbox}
\section[\texorpdfstring{\mr{मथळा} (mathḷā)}{मथळा (mathḷā)}]{\mr{मथळा} (mathḷā) — a headline}

\label{lesson:MR-C310-mathla}



\begin{quote}
Before the new one: say the Marathi for a daily paper, then the Marathi for a weekly.
\end{quote}

\subsection*{You'll want to know: \mr{मथळा}}
\textbf{\mr{मथळा}} — \emph{mathḷā} — \textquotedblleft{}a headline\textquotedblright{}.

\begin{quote}
\mr{आजच्या} \mr{वृत्तपत्राचा} \mr{मथळा} \mr{पावसाबद्दल} \mr{आहे}.
\end{quote}

\begin{quote}
\emph{ājacyā vṛttapatrācā mathḷā pāvsābaddal āhe.}
\end{quote}

\begin{quote}
Today's newspaper headline is about the rain.
\end{quote}

Masculine. Plural: \mr{मथळे}, \emph{mathḷe}.

Plural \mr{मथळे}.

\begin{cousinweb}
\textbf{\mr{मथळा}} (\emph{mathḷā}) is the Marathi word for a headline; its origin is not settled.
\end{cousinweb}

\begin{grammarlens}[title={What you've built}]
One more word for telling what happened.
\end{grammarlens}

\subsection*{Your turn}
\begin{itemize}
\raggedright
  \item \emph{Say it:} \emph{mathḷā}
  \item \emph{Say it:} \emph{mathḷā}, once more
  \item \emph{Say it:} \emph{ājacyā vṛttapatrācā mathḷā pāvsābaddal āhe.} — Today's newspaper headline is about the rain.
  \item \emph{Recall:} say the Marathi for a daily paper, then the Marathi for a weekly, then say \emph{mathḷā} again
\end{itemize}

\begin{tcolorbox}[breakable,colback=teal!4,colframe=teal!35!black,title={Before you move on}]
What does \mr{मथळा} mean? (\textquotedblleft{}A headline\textquotedblright{}.) Say it once more.
\end{tcolorbox}
\section[\texorpdfstring{\mr{छापणे} (chāpṇe)}{छापणे (chāpṇe)}]{\mr{छापणे} (chāpṇe) — to print}

\label{lesson:MR-C310-chapne}



\begin{quote}
Before the new one: say the Marathi for a weekly, then the Marathi for a headline.
\end{quote}

\subsection*{You'll want to know: \mr{छापणे}}
\textbf{\mr{छापणे}} — \emph{chāpṇe} — \textquotedblleft{}to print\textquotedblright{}.

\begin{quote}
\mr{त्या} \mr{छापखान्याने} \mr{हजार} \mr{पुस्तकं} \mr{छापली}.
\end{quote}

\begin{quote}
\emph{tyā chāpkhānyāne hajār pustakaṁ chāplī.}
\end{quote}

\begin{quote}
That press printed a thousand books.
\end{quote}

Past, with -\mr{ने} (he …): \mr{त्याने} \mr{छापलं}, \emph{tyāne chāplaṁ}.

What is printed agrees: \mr{पुस्तक} \mr{छापलं}, \mr{बातमी} \mr{छापली}.

\begin{cousinweb}
\textbf{\mr{छापणे}} (\emph{chāpṇe}) is the Marathi verb for printing and stamping; its origin is not settled.
\end{cousinweb}

\begin{grammarlens}[title={What you've built}]
One more word for telling what happened.
\end{grammarlens}

\subsection*{Your turn}
\begin{itemize}
\raggedright
  \item \emph{Say it:} \emph{chāpṇe}
  \item \emph{Say it:} \emph{chāpṇe}, once more
  \item \emph{Say it:} \emph{tyā chāpkhānyāne hajār pustakaṁ chāplī.} — That press printed a thousand books.
  \item \emph{Recall:} say the Marathi for a weekly, then the Marathi for a headline, then say \emph{chāpṇe} again
\end{itemize}

\begin{tcolorbox}[breakable,colback=teal!4,colframe=teal!35!black,title={Before you move on}]
What does \mr{छापणे} mean? (\textquotedblleft{}To print\textquotedblright{}.) Say it once more.
\end{tcolorbox}
